



L23_Amazone Sale_2026_Case Study_Coding Session.pdf
PDF
Example Lecture Scribe (structure only):

Topic: Random Variables and Expectations

Definition
A random variable maps outcomes of an experiment to real numbers.
Key Formula
Expectation:
��
[
��
]
=
∑
��
⋅
��
(
��
=
��
)
E[X]=∑x⋅P(X=x)
Explanation
Expectation represents the average outcome over many trials.
Worked Example
If 
��
X takes values {1,2,3} with probabilities {0.2, 0.5, 0.3}:
��
[
��
]
=
1
(
0.2
)
+
2
(
0.5
)
+
3
(
0.3
)
=
2.1
E[X]=1(0.2)+2(0.5)+3(0.3)=2.1
Insight
Expectation does not indicate the most likely value, but the long-run average.

Instruction:

Using the same structure, generate a lecture scribe for the provided lecture (Amazon Black Friday Sale 2026 case study on tail bounds and system reliability). The scribe should function as exam reference material and include all key definitions, formulas (Markov, Chebyshev, Chernoff), derivations, comparisons, and worked examples exactly as presented in the lecture. Keep the content structured, concise, and aligned with the lecture flow.

Here is your lecture scribe following the exact structure and aligned with the provided lecture 

:

Topic: Tail Bounds and System Reliability (Amazon Black Friday Sale 2026 Case Study)
Definition
Tail probability refers to the probability that a random variable deviates significantly from its mean:

P
(
X
≥
k
)
P(X≥k)
In system reliability, this represents the chance of extreme events (traffic spikes) that may cause system failure (SLA breach).

Key Formulas
1. Markov’s Inequality
P
(
X
≥
a
)
≤
E
[
X
]
a
P(X≥a)≤ 
a
E[X]
​
 
2. Chebyshev’s Inequality
P
(
∣
X
−
μ
∣
≥
a
)
≤
σ
2
a
2
P(∣X−μ∣≥a)≤ 
a 
2
 
σ 
2
 
​
 
3. Chernoff Bound
P
(
X
≥
a
)
≤
min
⁡
t
>
0
e
−
t
a
M
X
(
t
)
P(X≥a)≤ 
t>0
min
​
 e 
−ta
 M 
X
​
 (t)
where:

M
X
(
t
)
=
E
[
e
t
X
]
M 
X
​
 (t)=E[e 
tX
 ]
Explanation
Real-world systems (like Amazon servers) experience random traffic, often modeled using Poisson distribution 

 (page 3).

The mean load may appear safe, but rare spikes (tail events) can cause catastrophic failures (page 5: “The mean is safe. The tail is catastrophic.”).

Exact probabilities are hard to compute → we use upper bounds:

Markov → uses only mean

Chebyshev → uses mean + variance

Chernoff → uses full distribution (MGF)

Derivations
Markov Inequality
From non-negative random variable:

E
[
X
]
≥
a
⋅
P
(
X
≥
a
)
⇒
P
(
X
≥
a
)
≤
E
[
X
]
a
E[X]≥a⋅P(X≥a)⇒P(X≥a)≤ 
a
E[X]
​
 
Chebyshev Inequality
Using variance:

V
a
r
(
X
)
=
E
[
(
X
−
μ
)
2
]
Var(X)=E[(X−μ) 
2
 ]
Apply Markov on 
(
X
−
μ
)
2
(X−μ) 
2
 :

P
(
∣
X
−
μ
∣
≥
a
)
≤
V
a
r
(
X
)
a
2
P(∣X−μ∣≥a)≤ 
a 
2
 
Var(X)
​
 
Chernoff Bound
Using exponential transformation:

P
(
X
≥
a
)
=
P
(
e
t
X
≥
e
t
a
)
P(X≥a)=P(e 
tX
 ≥e 
ta
 )
Apply Markov:

P
(
X
≥
a
)
≤
E
[
e
t
X
]
e
t
a
=
e
−
t
a
M
X
(
t
)
P(X≥a)≤ 
e 
ta
 
E[e 
tX
 ]
​
 =e 
−ta
 M 
X
​
 (t)
Optimize over 
t
t for tightest bound.

Worked Example (Amazon Case Study)
Scenario
Traffic follows: Poisson(
λ
=
500
λ=500) 

 (page 11)

SLA failure threshold: 
k
=
600
k=600

Goal: Compute 
P
(
X
≥
600
)
P(X≥600)

Results (from lecture visualization) 

 (page 12)
Exact Probability ≈ 1.4%

Markov Bound ≈ 90.9%

Chebyshev Bound ≈ 20%

Chernoff Bound ≈ 8.8%

Interpretation
Markov → extremely loose (overestimates risk)

Chebyshev → better but still loose

Chernoff → closest to reality

Comparison
Method	Requires	Decay Rate	Tightness	Use Case
Markov	Mean (
μ
μ)	
1
/
k
1/k (Linear)	Very Loose	Quick rough estimate
Chebyshev	Mean + Variance	
1
/
k
2
1/k 
2
  (Polynomial)	Moderate	When variance known
Chernoff	Full MGF	
e
−
k
e 
−k
  (Exponential)	Tight	Real system design
(From decision matrix, page 16 

)

Insight
Better inputs → tighter bounds (page 6):

Mean → weak prediction

Variance → better

Full distribution → best

Engineering impact:

Overestimating probability ⇒ over-provisioning servers

Waste multiplier:

Waste
=
Bound Estimate
Exact Probability
Waste= 
Exact Probability
Bound Estimate
​
 
Example (page 13):

Markov → ~107,000× waste

Chebyshev → ~6,422× waste

Chernoff → ~10.7× waste

Final Takeaway
The goal is not exact probability, but safe and tight upper bounds.

Markov is safest but wasteful, Chernoff is precise but requires more information.

In high-stakes systems (like Black Friday), tight bounds directly save millions in infrastructure costs.

If you want, I can convert this into perfect PPT slides format or exam cheat sheet (1-page revision).






